\documentclass[10pt,conference,compsocconf]{IEEEtran}

\usepackage{amsmath}
\usepackage{booktabs}
\usepackage{graphicx}
\usepackage{hyperref}
\usepackage{url}

\hypersetup{hidelinks}

\begin{document}

\title{Predicting Coronary Heart Disease from BRFSS Survey Data}

\author{
  Ahmed Azaiez, David Gorgiev, and Peter Guggenbach\\
  \textit{CS-433 Machine Learning, EPFL}
}

\maketitle

\begin{abstract}
We predict whether a respondent reports myocardial infarction or coronary
heart disease (MICHD) from 321 BRFSS survey variables~\cite{cdc2015data}.  Only 8.8\% of the
328,135 training examples are positive, so we optimize F1 rather than
accuracy.  Our NumPy-only system combines questionnaire-aware preprocessing,
categorical and missing-value indicators, nonlinear features, logistic
regression, and a threshold fitted for F1.  Stratified five-fold
cross-validation improves from $0.3946\pm0.0039$ F1 for the MSE-gradient
baseline to $0.4437\pm0.0022$.  Training on all labelled rows yields 0.455 F1
and 0.882 accuracy on AIcrowd.  The largest gains come from the decision
threshold and preprocessing; regularization has little effect.  We also
identify a survey-design target proxy and discuss why the leaderboard result
must not be interpreted as prospective clinical performance.
\end{abstract}

\section{Problem and Data}

The task is binary classification with labels $y\in\{-1,+1\}$, where $+1$
denotes MICHD.  The inputs contain health, demographic, lifestyle, and survey
administration variables.  The strong class imbalance makes accuracy alone
misleading: predicting only the majority class would be accurate but would
detect no positive cases.  We therefore use positive-class F1,
$2TP/(2TP+FP+FN)$, as the model-selection metric and report accuracy as a
secondary measure.  Because the survey records existing diagnoses, our model
predicts reported MICHD status, not the future causal risk of developing it.

\section{Methods}

\subsection{Preprocessing and Features}

All fitted preprocessing is learned only from the current training fold and
then applied unchanged to its validation fold.  Questionnaire codes from the
2015 BRFSS instrument~\cite{cdc2015questionnaire} for
``none'', ``don't know'', and ``refused'' are converted to zero or missing as
appropriate.  Mixed units---for example kilograms versus pounds and weekly
versus monthly frequencies---are converted to common units, and deterministic
skip logic is resolved before imputation.  Remaining missing values use
training-fold medians, and every output column is standardized with
training-fold statistics.

Categorical variables marked by our 321-column feature registry are one-hot
encoded, with a separate unknown level.  Missing indicators preserve whether
a numeric value was imputed.  We add squared standardized age and BMI, while
14 heavy-tailed count variables are clipped at the training 99.5th percentile
and transformed with $\log(1+x)$.  Module-question flags retain information
from sparse optional survey sections.  Constant and more-than-95\%-missing raw
columns are otherwise removed.

Because a linear score cannot create interactions, we add 131 fixed columns:
age and sex interactions with risk factors, pairwise risk-factor products,
linear and squared counts of cardiovascular, chronic, and disability factors,
and interactions of general health with age, sex, and risk factors.  Seven
further columns interact whether the optional daily-aspirin question was
answered with age, sex, general health, high blood pressure, high cholesterol,
diabetes, and blood-pressure medication.  Together with the preceding
transformations, the selected design matrix has 729 columns before these
interactions and 736 columns in the final submitted design.

\subsection{Model, Threshold, and Tuning}

We implemented the six required regression methods from scratch using NumPy.
The selected classifier is logistic regression.  For labels converted to
$\{0,1\}$, it minimizes the mean binary cross-entropy
\begin{equation}
  \mathcal{L}(w)=\frac{1}{N}\sum_{i=1}^{N}
  \left[\log(1+e^{x_i^\top w})-y_i x_i^\top w\right]
\end{equation}
by full-batch gradient descent from $w=0$.  Scores $x^\top w$ are converted to
labels using the cutoff that maximizes F1 on the training rows.  This cutoff is
refitted inside every fold, so validation labels are never used to choose it.

We searched step sizes $\gamma\in\{0.05,0.08,0.10,0.12\}$ at 500, 1000,
2000, and 3000 iterations.  At the selected optimizer setting we evaluated
$\lambda\in\{0,10^{-6},10^{-5},10^{-4},10^{-3},10^{-2}\}$.  To avoid paying
for negligible differences, we chose the cheapest and simplest setting
within 0.0005 F1 of the best mean.  This selected $\gamma=0.10$, 1000
iterations, and $\lambda=0$.

\subsection{Evaluation Protocol}

Every comparison uses the same five stratified folds and seed 42 on all
328,135 labelled rows.  Each fold contains approximately 262,508 training and
65,627 validation rows.  Imputation, scaling, clipping, encodings, model
weights, and the threshold are fitted anew on its training partition.  We
report mean and population standard deviation across folds.  The final
AIcrowd file is produced by fitting the selected pipeline once on all labelled
rows and predicting the 109,379 test rows.

\section{Results}

Table~\ref{tab:ablation} shows a controlled chain.  The first three rows give
reference configurations: squared loss, logistic loss, and an F1-specific
cutoff.  The largest single gain is choosing the decision threshold for F1.
The codebook-aware feature pipeline contributes a further 0.0134 F1, and
optimizer tuning improves all five folds by another 0.0041.  Engineered
features add 0.0035 F1 in all five folds; aspirin interactions add 0.0007 in
four of five folds, within the fold spread.  We retain them because this frozen
configuration gave the best AIcrowd score.  The complete system is 0.0490 F1
above the original baseline.

\begin{table}[t]
  \centering
  \caption{Five-fold ablation on all training rows.  $\Delta$ is relative to
  the preceding row, except the L2 and no-flags comparisons, which use the
  tuned optimizer row.}
  \label{tab:ablation}
  \small
  \setlength{\tabcolsep}{3pt}
  \begin{tabular}{lcc}
    \toprule
    Configuration & CV F1 & $\Delta$ \\
    \midrule
    MSE GD + prevalence cutoff & $0.3946\pm.0039$ & --- \\
    Logistic + prevalence cutoff & $0.4023\pm.0029$ & +.0076 \\
    Logistic + F1 cutoff & $0.4219\pm.0022$ & +.0197 \\
    + questionnaire-aware features & $0.4353\pm.0022$ & +.0134 \\
    + tuned optimizer & $0.4394\pm.0036$ & +.0041 \\
    + L2 ($\lambda=10^{-3}$) & $0.4396\pm.0032$ & +.0002 \\
    + engineered features & $0.4430\pm.0026$ & +.0035 \\
    + aspirin interactions & $\mathbf{0.4437\pm.0022}$ & +.0007 \\
    \addlinespace
    Tuned model, no module flags & $0.4295\pm.0037$ & $-.0099$ \\
    \bottomrule
  \end{tabular}
\end{table}

Convergence explains the last gain.  At $\gamma=0.05$, increasing from 500 to
2000 iterations raises F1 from 0.4353 to 0.4393, after which improvement is
negligible.  Doubling the step size reaches the same region in 1000 iterations.
The largest regularized mean, 0.4396 at $\lambda=10^{-3}$, differs from the
unregularized 0.4394 by only 0.0002; $\lambda=10^{-2}$ reduces F1 to 0.4369.
We therefore retain the simpler unregularized model.

\begin{table}[t]
  \centering
  \caption{Mean five-fold F1 across the optimizer grid.}
  \label{tab:tuning}
  \scriptsize
  \setlength{\tabcolsep}{3pt}
  \begin{tabular}{ccccc}
    \toprule
    $\gamma\backslash$ iterations & 500 & 1000 & 2000 & 3000 \\
    \midrule
    .05 & .4353 & .4375 & .4393 & .4395 \\
    .08 & .4371 & .4386 & .4391 & .4391 \\
    .10 & .4374 & \textbf{.4394} & .4391 & .4388 \\
    .12 & .4377 & .4393 & .4391 & .4387 \\
    \bottomrule
  \end{tabular}
\end{table}

The initial uploaded logistic model scored 0.418 F1 and 0.897 accuracy on
AIcrowd.  The final system predicts 12.90\% of test rows as positive and scores
\textbf{0.455 F1 with 0.882 accuracy}.  The lower accuracy is expected: an
F1-oriented cutoff predicts more positives, trading some majority-class
accuracy for better detection of the minority class.  The final AIcrowd F1 is
0.0113 above its CV estimate, which is compatible with sampling variation and
differences between the public test set and our folds.

\section{Limitations and Discussion}

Threshold selection and semantically correct preprocessing matter more here
than regularization.  However, one useful feature is also a serious limitation.
\texttt{HAREHAB1} asks about rehabilitation \emph{following a heart attack} and
is only shown to respondents reporting one; all 773 training respondents for
whom it is answered are positive.  A separate ablation estimated that sparse
module flags, dominated by this variable, add about 0.0099 F1.  The feature is
present in both provided train and test data, but it is a direct target proxy,
a recognized form of data leakage~\cite{kaufman2012leakage}.  We retain the
module flags and aspirin interactions to reproduce the benchmark submission;
removing the module flags lowers CV F1 from 0.4394 to 0.4295.  Consequently,
this system is suitable for the defined benchmark, not for prospective
screening.  A clinically meaningful
follow-up should remove all post-outcome questions and report the resulting
performance separately.

The same folds were used for hyperparameter selection and the final CV number,
so the reported mean may contain modest selection optimism.  The independent
AIcrowd score provides an additional check, but repeated leaderboard queries
must not replace local validation.  Nested CV or a sealed holdout would give a
cleaner final generalization estimate.

\section{Conclusion and Reproducibility}

A transparent NumPy pipeline combining logistic loss, an F1-specific cutoff,
and questionnaire-aware feature engineering reaches 0.455 AIcrowd F1.  The
ablation shows that the threshold and feature semantics account for most of
the improvement, while optimizer tuning provides a smaller consistent gain
and regularization is unnecessary at the tested scale.

The hand-in repository is
\url{https://github.com/Guggi99/ml-project1-2026-submission}.  Experiments share
one cross-validation pipeline and store their folds, seed, transformations,
model, and cutoff rule.  With the three course CSV files in \texttt{data/},
\texttt{python run.py --data-dir data --submission submission.csv}
reproduces the frozen submission (SHA-256 prefix \texttt{0a97469b}), which
scored 0.455 F1 and 0.882 accuracy on AIcrowd versus a CV estimate of 0.4437.
Every push runs unit, formatting, and course grading tests.  AI assistants were
used in this project in compliance
with EPFL rules.  All design decisions, code and text were reviewed and are the
responsibility of the authors.

\begin{thebibliography}{1}

\bibitem{cdc2015data}
Centers for Disease Control and Prevention, ``2015 BRFSS survey data and
documentation,'' 2015. [Online]. Available:
\url{https://www.cdc.gov/brfss/annual_data/annual_2015.html}

\bibitem{cdc2015questionnaire}
Centers for Disease Control and Prevention, ``2015 BRFSS questionnaire,''
2015. [Online]. Available: \url{https://www.cdc.gov/brfss/questionnaires/}

\bibitem{kaufman2012leakage}
S.~Kaufman, S.~Rosset, C.~Perlich, and O.~Stitelman, ``Leakage in data mining:
Formulation, detection, and avoidance,'' \emph{ACM Trans. Knowl. Discov. Data},
vol.~6, no.~4, pp.~15:1--15:21, 2012, doi: 10.1145/2382577.2382579.

\end{thebibliography}

\end{document}
